% SPDX-License-Identifier: Apache-2.0
\section{The Last Thing That Happened}
\label{sec:x-tracer}

A design on a board says almost nothing. Four LEDs and a serial line
are what a program has, and neither answers the question a post-mortem
asks, which is what the machine did in the cycles before it stopped.

The answer is usually already computed. The Vreteno core drives the
instruction it retired and the register it wrote, every cycle, and on
the board those wires go nowhere. \code{Tracer<N>} takes a word a
cycle and keeps the last \code{N} of them in a ring, so that what is
held is always the most recent window; a host reads the window out
over AXI-Lite, oldest first.

Freezing is the point of the \code{halt} input. A window that keeps
filling after the machine stops is a window of whatever happened next,
which is nothing anybody is asking about, so with \code{CTRL\_FREEZE}
set the first cycle of \code{halt} clears the run bit. The buffer then
holds still, and the bit a host reads back says why it stopped.

An entry is 128 bits because a retirement is: a program counter, an
instruction, the register written and the word written come to a
hundred and one bits, and a width that is a whole number of bus words
is one the readout can walk without shifting. The four words of an
entry are read at \code{0x10} to \code{0x1c}, and the cursor at
\code{0x08} says which held entry they belong to.

The run offers twelve entries into a ring of eight, raises
\code{halt}, and offers four more that the buffer must refuse. The
readout that follows is what a post-mortem wants: the last eight
things the machine did, oldest first, and nothing after the halt.

\begin{figure}[htbp]
\centering
\input{tracer_timing.tex}
\caption{The buffer filling: \code{head} walks the ring, \code{count}
stops at eight, and \code{ctrl} loses its run bit when \code{halt}
rises.}
\label{fig:x-tracer}
\end{figure}

\lstinputlisting[caption={\code{ex\_tracer.rs}. Run with \code{bazel run //lib/examples:ex\_tracer}.},label={lst:x-tracer}]{lib/examples/ex_tracer.rs}

\lstinputlisting[style=ascii,caption={What \code{ex\_tracer} printed when the build ran it: eight held, the run bit cleared by the halt, and the window read out oldest first.},label={lst:x-tracer-out}]{out_tracer.txt}
